% HD-S47. GPT-6 Astra Pro /root/sixfolds, inherited effort; 21 November 2026.
\section{HD-S47: two-dimensional blocks for rank-three generators}
\label{sec:two-block-rank-three-scope}

Let \(K\subset\mathbf C\) have characteristic zero, and let
\(N\subset GL_{2r}(K)\), \(r\ge8\), be finite and absolutely
irreducible. Suppose \(N\) is generated by its actual moving-rank-three
elements, has no nonidentity element of moving rank one or two,
and has a complex system of \(r\) two-dimensional blocks.
Then the kernel \(T\) on these blocks is characteristic in
\(N\), and the original block projectors have entries in \(K\).
Every actual \(K\)-linear overgroup normalizing \(N\) preserves
these descended blocks.

For every imaginary quadratic field, a natural primitive
sixfold source with such a normal subgroup therefore has
period rank at most seven. This removes the two-dimensional
block alternative from form3 of HD-R15. The argument does not
assume that the extension of the block permutation group splits.

\paragraph{1. The exact outside generators.}
Write \(V_{\mathbf C}=\bigoplus_{i=1}^rV_i\), with
\(\dim V_i=2\). The fixed-space calculation of HD-R15
identifies the fixed vectors on a cycle of blocks with the
kernel of its return map. A rank-three generator outside
\(T\) must therefore exchange exactly one pair \(i,j\).
The codimension on that pair is
\(2+\operatorname{rank}(Q-I)\), where \(Q\) is its
two-dimensional return map. If the second term were one,
the generator would be the identity on all other blocks,
and its square would have moving rank two: its two diagonal
return maps are conjugate to \(Q\), each of moving rank
one. This is forbidden. Thus \(Q=I\). The remaining
codimension is one, on a single third block \(k\). The
square is now supported only on that block; it would have
moving rank one unless the nonidentity eigenvalue there
were \(-1\). Hence every outside generator is literally
\[
 x|_{V_i\oplus V_j}=\begin{pmatrix}0&A\\A^{-1}&0\end{pmatrix},
 \qquad x|_{V_k}=R,\quad R^2=I,\quad\operatorname{rank}(R-I)=1,
 \qquad x|_{V_\ell}=I\ (\ell\notin\{i,j,k\}).
 \tag{S47.1}
\]
In particular it is an involution. The graph of its exchanged
pairs is connected; otherwise a connected-component sum is
invariant under all generators. Transpositions along a connected
graph generate the full symmetric group. Consequently
\[
                           N/T\simeq S_r
 \tag{S47.2}
\]
by the actual block action. No pure-transposition lifts have
been asserted.

Choose \(n_i\in N\) carrying \(V_1\) to \(V_i\), and
use its restriction to transport a basis of \(V_1\).
Normality of \(T\) then identifies all its coordinate images
with one finite group \(L\subset GL_2(\mathbf C)\), and
the literal coordinate restriction gives a subdirect inclusion
\(T\subset L^r\). Let \(H_i\) denote the stabilizer of
\(V_i\) in \(N\). Its representation on \(V_i\) is
irreducible: a proper invariant subspace would induce a
proper invariant direct sum under the map
\(n\otimes v\mapsto\rho(n)v\). Also
\(H_i/T\simeq S_{r-1}\) by (S47.2).

\paragraph{2. The coordinate group is irreducible, including the scalar case.}
We first exclude scalar \(L\). Since \(T\) is scalar
on the specified block, the actual projective map
\(H_i\to PGL(V_i)\) kills \(T\), and therefore factors
through the literal quotient \(S_{r-1}\). For \(r-1\ge7\),
its finite image must be trivial or \(C_2\). Indeed the
normal subgroups of \(S_{r-1}\) give only these quotients
or the faithful image \(S_{r-1}\), and the latter is not
in the finite \(PGL_2(\mathbf C)\) list
\(C_a,D_a,A_4,S_4,A_5\). A linear group with cyclic
projective image is generated by its scalar subgroup and
one lift, and so has invariant eigenlines. This contradicts
the irreducibility of \(H_i\) on \(V_i\).
This is a projective quotient argument that includes every
possible scalar central extension or spin lift; no splitting
of such an extension is assumed.

Now suppose \(L\) is reducible but nonscalar. Complete
reducibility for finite groups makes it diagonal on two
distinct characters, and \(T\subset L^r\) is abelian.
Let \(\mathcal S_i\) be the two distinct \(T\)-characters
on \(V_i\). Conjugation by \(N\) permutes these pairs
according to the full \(S_r\). Hence
\(|\mathcal S_i\cap\mathcal S_j|\) is constant for
distinct \(i,j\), and is zero, one or two.

If it is two, there are exactly two global \(T\)-character
spaces, each of dimension \(r\). Irreducibility makes
\(N\) exchange them, so one of its rank-three generators
exchanges them. Such an exchange has moving rank at least
\(r>3\), a contradiction.
If the intersection size is one, the distinct two-element
sets form a star when there are at least four of them:
from \(\{a,b\}\) and \(\{a,c\}\), a pair omitting
\(a\) must be \(\{b,c\}\); a fourth distinct pair
cannot meet all three. Thus all the pairs have a unique
common character. It is fixed by \(N\), and its proper
nonzero character space contradicts irreducibility.

Finally suppose the pairs are disjoint. Their \(2r\)
character spaces are distinct lines, permuted by \(N\).
The kernel of this line permutation is exactly \(T\):
\(T\) is diagonal on them, while fixing all these lines
fixes every original block. Thus the line permutation
group is the same actual quotient \(N/T=S_r\).
For each block, \(H_i/T=S_{r-1}\) must exchange its
two character lines, since otherwise \(H_i\) would be
reducible. Its map to \(S_2\) is therefore the sign map.
An outside generator projects to a transposition \((i\ j)\).
For every \(k\ne i,j\), this is a transposition in the
point stabilizer \(S_{r-1}\), so it exchanges the two
character lines of \(V_k\). Each such fixed block then
contributes at least one to its moving rank. The exchanged
pair contributes at least two, giving total moving rank
at least \(2+(r-2)=r>3\), again a contradiction.
This exhausts the reducible case, and proves that \(L\)
is irreducible.

The finite projective classification used here and below is
Nguyen--van der Put--Top, \emph{Algebraic subgroups of
\(GL_2(\mathbf C)\)}, Indagationes Mathematicae19 (2008),
287--297, Theorem2, printed p.288,
\url{https://doi.org/10.1016/S0019-3577(08)80004-3};
the author-hosted original is
\url{https://www.math.rug.nl/~top/INDAG171.pdf}.
That theorem explicitly attributes the finite projective
list to Klein. Only this complete list is used here.

\paragraph{3. The intrinsic kernel without a splitting assumption.}
If \(L\) is solvable, then \(T\subset L^r\) is
solvable. Every solvable normal subgroup of \(N\) maps
trivially to \(S_r\), so
\(T=\operatorname{Rad}(N)\) is characteristic.

In the nonsolvable case, the literal projective image is
\(L/Z=A_5\), where \(Z=L\cap\mathbf C^*I_2\).
Coordinate projectivization gives
\[
 P:=\operatorname{im}(T\longrightarrow A_5^r),\qquad
 Z_0:=\ker(T\longrightarrow P)=T\cap Z^r.
 \tag{S47.3}
\]
The group \(P\) is subdirect. A subdirect subgroup of a
power of a nonabelian simple group is a product of twisted
diagonal subgroups on a partition of its coordinates.
Here is the actual-map proof. Induct on the number of
coordinates. Projection to the first coordinates has,
by induction, a product of simple diagonal factors as
image. The kernel in the last factor is normal and hence
either the full last simple group or trivial. In the first
case the new factor is independent. In the second case
the subgroup is the graph of a surjection from that product
to the last simple group. The images of the simple factors
are commuting normal subgroups, so exactly one factor
maps isomorphically onto the last group and all the others
map trivially. This links the last coordinate to precisely
one existing diagonal factor, through the displayed
isomorphism. The induction supplies all the claimed maps.

Conjugation by \(N\) permutes coordinates through \(S_r\)
and conjugates their projective matrix groups. It preserves
the equivalence relation of belonging to one diagonal factor.
Since \(S_r\) is two-transitive, that relation is either
equality or the universal relation. Thus \(P=A_5^r\),
or \(P\) is a single twisted diagonal \(A_5\), with
each coordinate projection injective.

The latter case is impossible by the exact generator support.
Take an outside generator as in (S47.1), and choose a block
\(\ell\notin\{i,j,k\}\). It is the identity on that
block. Its induced automorphism of \(P\) therefore has
unchanged \(\ell\)-th coordinate, and the injectivity
of that projection makes the automorphism the identity.
In coordinate \(k\), this says that the nontrivial
projective involution defined by \(R\) centralizes
\(A_5\subset PGL_2(\mathbf C)\). This cannot happen:
the two fixed points of that involution would be preserved
by \(A_5\); its map to their permutation group is trivial
because \(A_5\) is perfect, so it would fix both points
individually and lie in the abelian diagonal torus.
Consequently
\[
          P=A_5^r,\qquad T/Z_0=A_5^r.
 \tag{S47.4}
\]

The scalar-coordinate subgroup \(Z_0\) is abelian and
normal in \(N\). Any solvable normal subgroup maps
trivially to \(S_r\) and then trivially to \(A_5^r\).
Thus \(\operatorname{Rad}(N)=Z_0\). In the actual
quotient \(\overline N=N/Z_0\), the specified base
\(T/Z_0=A_5^r\) is minimal normal, since the top \(S_r\)
permutes its simple factors transitively. Its centralizer
in \(\overline N\) is trivial. Indeed a centralizing
element must have trivial coordinate permutation, so lies
in the base, whose centre is trivial. Any other minimal
normal subgroup would intersect the base trivially and
centralize it. Thus this base is the unique minimal normal
subgroup, and
\[
 T=\text{the full preimage of the unique minimal normal
 subgroup of }N/\operatorname{Rad}(N).
 \tag{S47.5}
\]
This proves characteristicness. The quotient used in this
argument is the literal one induced by coordinate
projectivization and the block action. No isomorphism
to a split wreath product is required or asserted.

\paragraph{4. Distinct kernel types and actual field descent.}
Each block is now an irreducible \(T\)-module. Equality
of its isomorphism type with another block is an equivalence
relation invariant under the full \(S_r\). It is therefore
either equality or the universal relation. The universal
case would give an actual fixed complex tensor structure,
which is excluded by HD-R14. Explicitly, let \(W\) be
one two-dimensional irreducible type and set
\(M=\operatorname{Hom}_T(W,V_{\mathbf C})\). The evaluation
\[
       W\otimes_{\mathbf C}M\xrightarrow{\sim}V_{\mathbf C},
                         w\otimes f\longmapsto f(w)
 \tag{S47.6}
\]
is an isomorphism: choose one intertwiner onto each block
to obtain its explicit direct-sum inverse. Matrix-coefficient
averaging realizes \(\mathbf C\rho(T)\) as
\(\operatorname{End}(W)\otimes I_M\). Conjugation by
each \(n\in N\) preserves this algebra. Its matrix-algebra
automorphism is inner: from the images of the standard
matrix units, choose \(v\ne0\) in the image of the first
diagonal unit and use the columns \(v,E'_{21}v\) to
construct an implementing invertible matrix \(C_n\).
Then \((C_n^{-1}\otimes I)\rho(n)\) commutes with the
matrix algebra and equals \(I\otimes D_n\). Thus the
literal action is \(C_n\otimes D_n\) on the fixed
\(2\)-by-\(r\) tensor decomposition. HD-R14 excludes
this for the absolutely irreducible rank-three-generated
group of dimension \(2r\ge16\).

The block \(T\)-modules are consequently pairwise
inequivalent. Their actual character projectors are
\[
 p_i=\frac2{|T|}\sum_{t\in T}\chi_i(t^{-1})\rho(t).
 \tag{S47.7}
\]
As in (S46.11), averaging on intertwiner spaces proves
that these are exactly the original block projections,
and the complete primitive central projectors of the
image algebra. They lie in \(E\rho(T)\) for
\(E=K(\zeta_{\exp T})\). Its Galois group permutes
them and commutes with the natural \(S_r\)-action,
because all \(N\)-matrices have entries in \(K\).
The centralizer of that permutation action is trivial.
Therefore
\[
       p_i\in M_{2r}(K),\qquad U_i=p_iK^{2r},\qquad
       U_i\otimes_K\mathbf C=V_i.
 \tag{S47.8}
\]
Characteristicness of \(T\) makes every actual normalizer
of \(N\) permute these same central projectors. This
proves the asserted descent for the original blocks.

\paragraph{5. The original sixfold factor and its rank-two quotient.}
Let \(G\) be the full original coefficient group, containing
and normalizing \(N\), with exact imaginary quadratic field
\(K\). Its prescribed invariant positive Hermitian form
makes the inequivalent \(T\)-blocks orthogonal. Let \(H\)
stabilize \(U=U_1\), let \(\sigma\) be its absolutely
irreducible action there, and put \(H_0=\ker\sigma\).
The actual induction map is \(g\otimes u\mapsto\rho(g)u\);
the induced functions and inverse are
\[
 F_v(x)=p_1\rho(x)v,\quad F_v(hx)=\sigma(h)F_v(x),
 \qquad F\longmapsto\sum_{H\backslash G}\rho(x)^{-1}F(x).
 \tag{S47.9}
\]
Right translation gives the original coefficient action,
with inversion label \(Ht^{-1}\leftrightarrow tH\).
The inertia-cycle return-map identification gives these
same inverse maps on the ramified invariant stalks.
Induction of characters proves that the faithful rank-two
image retains exact character field \(K\).

Use the actual quotient curves
\(Y=\widetilde C/H\), \(C'=\widetilde C/H_0\), and
\(\alpha:\widetilde C\to C'\). For \(0\ne v\in U(K)\)
and \(q=h(v,v)\), the corner equality is literally
\[
 \frac{2(2r)}{|G|q}\sum_{g\in G}
  \operatorname{Re}h(v,\rho(g^{-1})v)g
 =\frac4{|H|q}\sum_{h_0\in H}
  \operatorname{Re}h(v,\sigma(h_0^{-1})v)h_0.
 \tag{S47.10}
\]
Off-\(H\) coefficients vanish and \(|G|=r|H|\).
After averaging the kernel, pullback carries the primitive
factor of \(JC'\) isogenously to this same abelian image,
with
\[
 \alpha_*\alpha^*=[|H_0|],\qquad
    (\alpha^*)^*\theta_{J\widetilde C}=|H_0|\theta_{JC'}.
 \tag{S47.11}
\]
The real Schur factors agree:
\(2|G|/((2r)q)=2|H|/(2q)\). For an originally specified
primitive vector \(w\), use (S46.17) with the coefficient
dimension \(2r\), namely the integer-cleared Fourier
operator of \(x\mapsto v h(w,x)/h(w,w)\).
Its Rosati multiplier is \(M^2h(v,v)/h(w,w)\).
These literal morphisms retain the original corner,
field action and polarized period differential after
finite parameter base change fixing the finite labels.

Finally the rank-two parabolic system on \(Y\) has
\(K\)-dimension six in degree one. If \(g=g(Y)\)
and \(b\) counts its true coefficient branch values,
\[
                    6=2(2g-2)+\sum_{j=1}^b c_j,
                          \qquad c_j\ge1.
 \tag{S47.12}
\]
For \(g=0\) this gives \(b\le10\), so at most seven
parameters. For \(g=1\) it gives at most six, and for
\(g=2\) at most \(3+2=5\); larger genus is impossible.
Finite monodromy data varies discretely over the fixed
pointed base, so these are bounds on the actual quotient
family's period rank. Equations (S47.9)--(S47.11) give
the same bound for the prescribed original sixfold.
Thus the two-dimensional block alternative has period
rank at most seven, as asserted.
